% !TEX root = ../nc-quadric-avatar.tex
% 2026-11-23 / 相対square完成と実際のtails認識。
\subsection{有限square族への限定的な拡張}
\label{subsec:qa-finite-square-family-opposite}

以下で基底は本稿の局所有限型 $k$-scheme の site に属するものとする。
次の補題は任意の四頂点代数を扱うのではなく、全幾何 fiber が既に分類された
square 型である有限局所自由族だけを扱う。

\begin{lemma}[分類済みfiberを持つ有限square族]
\label{lem:qa-classified-fiber-finite-square-recognition}
$S$ 上の有限局所自由代数 $B$ が四つの直交冪等元
$e_0,e_1,e_2,e_3$ を持ち、その右射影 $P_i=e_iB$ の Hom 束が
\[
 \left(\operatorname{rk}\operatorname{Hom}(P_i,P_j)\right)_{i,j}
 =\begin{pmatrix}1&2&2&4\\0&1&0&2\\0&0&1&2\\0&0&0&1\end{pmatrix}
\]
であるとする。対角成分は単位によって $\mathcal O_S$ と同定され、
逆向き成分は零とする。
さらに各幾何 fiber が、本稿で分類した square AS 型の初期四項の
有限 End 代数と、頂点と全積を保って同型であると仮定する。
古典型も許す。

このとき $B$ の実際の相対 perfect 圏から、$P_0$ を構造対象とする
相対 $\mathfrak H_2$ 族が構成される。四つの $P_i$ はこの同じ heart の
基底平坦対象となり、共通の fppf 局所 chart で compact である。
通常の canonical Serre 完成
\[
 C(B)=\bigoplus_{n\geq0}
 H^0\left(\Theta_B^{\otimes_B^{\mathbf L}n}\right),\qquad
 \Theta_B=\mathbf R\operatorname{Hom}_{B^e}(B,B^e)[2]
\]
は、各局所 affine chart とその許された基底変更上で右 Noetherian であり、
その tails はこの相対 heart と同値になる。
\end{lemma}
\begin{proof}
有限 directed 代数の normalized bar 解消により、$B$ は相対 smooth proper であり、
全 perfect Hom と相対 inverse Serre 関手は任意基底変更と可換する。
$H_{ij}=\operatorname{Hom}(P_i,P_j)$ とし、実際の cone
\[
 Z_1=\operatorname{Cone}(P_1\to H_{13}^{\vee}\otimes P_3),\qquad
 N_F=(P_0,P_2,P_3,Z_1)
\]
を取る。幾何 fiber の分類によって、二つの随伴
\[
 H_{01}\longrightarrow H_{13}^{\vee}\otimes H_{03},\qquad
 H_{02}\longrightarrow H_{23}^{\vee}\otimes H_{03}
\]
は各 fiber で単射である。曲線型では pencil の大域生成性、古典型では
通常の $\mathbb P^1\times\mathbb P^1$ の評価列による。
従って両者は基底上でも部分束の包含となる。

定理~\ref{thm:qa-relative-square-e0-e3-h2} の有限 cone の計算は、
これらの Hom 束と全積しか使わない。
したがって $N_F$ は相対 full strong exceptional 四項であり、
通常 End は階数二十四、$U_0$ の全射・部分束条件を満たす。
反対の二中間頂点を用いる窓 $N_G$ とは、B の相対 ruling 対合で交換される。

各幾何 fiber で $N_F$ が真正な normal foundation であることも、
同定理の体上の証明をそのまま適用できる。
cubic 開上では cubic 逆構成、一重壁では相対 ruling 交換の体上の形、
二重壁では既存の noncubic 認識を使う。
古典型では例えば
$(\mathcal O,\mathcal O(0,1),\mathcal O(1,1),\mathcal O(1,2))$
という normal 四項に一致する。
よって実際の射 $S\to\mathfrak F_0$ の全幾何 fiber が
$\mathfrak F_{\HH}$ に入り、既存の開埋込 $\Xi$ の逆を適用できる。
正規窓およびもう一方の ruling 窓から、四つの $P_i$ が
同じ相対 heart の基底平坦対象であることを得る。

全幾何 fiber の square helix の forward Ext 消滅と perfect-fiber 判定により、
$\mathbf R\operatorname{Hom}(P_i,t^mP_j)$ は $m\geq1$ に対して
次数零に集中する。とくに相対 AZ 命題の forward 条件が成立する。
従って $\bigoplus_iP_i$ は相対 AZ ample であり、
定理~\ref{thm:qa-relative-square-completion-tails} の証明を適用して
右 Noether 性と tails の認識を得る。
ここでは幾何 fiber の分類を入力とし、大域的な曲線 avatar の復元や
まだ得ていない左 Noether 性は使っていない。
\end{proof}

\subsection{反対代数の四頂点と添字の規約}

\begin{lemma}[square有限基礎の反対代数]
\label{lem:qa-square-foundation-opposite}
前補題の $B$ に対し $B^{\rm op}$ の標準冪等元を
\[
 f_i=e_{3-i}\qquad(0\leq i\leq3)
\]
と番号付け直す。このとき $B^{\rm op}$ も前補題の仮定を満たす。
\end{lemma}
\begin{proof}
$Q_i=f_iB^{\rm op}$ とすると
\[
 \operatorname{Hom}_{B^{\rm op}}(Q_i,Q_j)
 =f_jB^{\rm op}f_i
 =e_{3-i}Be_{3-j}
 =\operatorname{Hom}_B(P_{3-j},P_{3-i}).
\]
したがって Hom の階数表は再び square 型となり、積は元の積を逆転したものである。

幾何 fiber が曲線型 avatar $(D;P,F,G)$ に対応する場合を直接確認する。
\[
 P'=P^{-1},\qquad F'=FP^{-1},\qquad G'=GP^{-1}
\]
と置く。既存の曲線消滅補題により $F',G'$ は大域生成の次数二 pencil、
$F'G'=FGP^{-2}$ は very ample である。
また $G'(F')^{-1}=GF^{-1}$ は acyclic なので、二 pencil は同じ $D$ の
$(2,2)$ 埋込みを与える。$P'\ne\mathcal O_D$ も保たれる。
従って $(D;P',F',G')$ は再び square avatar である。

元の四線束を $(L_0,L_1,L_2,L_3)=(P,F,G,FG)$ とすれば、
新しい四線束は
\[
 (P',F',G',F'G')
 =(FGP^{-1})\otimes
 (L_3^{-1},L_2^{-1},L_1^{-1},L_0^{-1}).
\]
可逆層の双対は射の合成を逆転し、共通 tensor は Hom を保つ。
よって新しい directed End は、全積を含めて $B^{\rm op}$ と同型である。
本稿の avatar から AS 代数への逆構成により、これは分類済み square AS 型の
有限基礎となる。古典型でも、逆順双対 collection に
$\mathcal O(1,1)$ を共通 tensor すると標準 square collection を得る。
従って全幾何 fiber で所要の仮定が成立する。
\end{proof}

\begin{remark}[全整数AS代数で書く場合の反転]
本稿の正向きの規約が $A_{ji}=\operatorname{Hom}(E_i,E_j)$ であるとき、
上の有限基礎に対応する indexed opposite は
\[
 A^{\dagger}_{ji}=A_{3-i,\,3-j}
\]
であり、積を逆にする。
$i\mapsto-i$ だけを使うと初期位相がずれるので、同じ四位相表を主張するときは
定数 $3$ を保持する。
右単純の分解を双対化すると、元の添字 $i$ の右単純から
添字 $i-4$ の左単純の分解を得る。
新しい添字 $j$ に対して元の $i=7-j$ であり、距離 $k$ は $4-k$ に変わる。
従って $Q_{1,\rm new}$ は $Q_{2,\rm old}$、$Q_{2,\rm new}$ は
$Q_{1,\rm old}$ から得られる。
この変換は定義~\ref{def:qa-as-square-rotated-minimal} の square の四行を
それぞれ同じ四行へ送る。
ただし前補題の証明は、この全整数の双対性を未証明の入力にせず、
各体上の双対 avatar の有限 End を直接用いた。
\end{remark}

\subsection{inverse dualizing complexとopposite}

\begin{lemma}[Serre完成の反対代数への適合性]
\label{lem:qa-serre-completion-opposite}
$B$ を基底上有限局所自由で相対 smooth な代数とする。
$M^{\rm op}$ は $B$-双加群 $M$ の左右を交換した $B^{\rm op}$-双加群を表す。
有限対角解消に対して自然な同型
\[
 \Theta_{B^{\rm op}}\simeq(\Theta_B)^{\rm op}
\]
があり、非負 tensor 冪とその積を合わせると
\[
 \Pi_3(B^{\rm op})\simeq\Pi_3(B)^{\rm op}
\]
という内部次数付き dg 代数の自然な quasi-isomorphism を得る。
したがって通常の次数零部分についても
\[
 C(B^{\rm op})\simeq C(B)^{\rm op}
\]
が成立する。全ての比較は任意基底変更と可換する。
\end{lemma}
\begin{proof}
二つの enveloping 代数の間の因子交換
\[
 (B^{\rm op})^e=B^{\rm op}\otimes B
   \xrightarrow{\sim}B\otimes B^{\rm op}=B^e
\]
は対角双加群と有限局所射影対角解消を対応させる。
この解消を enveloping 代数で双対化し、標準の左右交換によって再び
双加群と解釈すれば、inverse dualizing complex の表示した同型を得る。
cohomological shift $[2]$ はこの比較と可換する。

左右交換は tensor 積について反モノイド的であり、
\[
 (M\otimes_B^{\mathbf L}N)^{\rm op}
 \xrightarrow{\sim}
 N^{\rm op}\otimes_{B^{\rm op}}^{\mathbf L}M^{\rm op}
\]
を与える。複体で書けば交換には cohomological Koszul 符号を付ける。
これを各 tensor 冪へ反復適用すると、積の順序を逆にした
内部次数付き dg 代数の quasi-isomorphism となる。
次数零 cohomology 上ではこの符号は通常の反対代数の積を与え、
内部次数に新たな符号を挿入しない。
有限局所射影解消、双対、tensor と因子交換はいずれも基底変更に自然なので、
全ての同型は任意基底変更と可換する。
\end{proof}

\begin{corollary}[square完成の左右Noether性]
\label{cor:qa-relative-square-completion-two-sided-noetherian}
実際の相対 curve-type square avatar から構成した $\Lambda$ について、
$C(\Lambda)$ は各局所 affine chart とその許された基底変更上で
左右 Noetherian である。
\end{corollary}
\begin{proof}
右 Noether 性は既存の tails 認識による。
反転した $\Lambda^{\rm op}$ は
補題~\ref{lem:qa-square-foundation-opposite} により
全幾何 fiber が分類済み square 型である有限局所自由族なので、
補題~\ref{lem:qa-classified-fiber-finite-square-recognition} により
$C(\Lambda^{\rm op})$ は右 Noetherian である。
前補題の同型によりこれは $C(\Lambda)^{\rm op}$ の右 Noether 性であり、
$C(\Lambda)$ の左 Noether 性を意味する。
必要な局所主張は中央の faithfully flat 基底変更によって降下する。
一般の任意 Noetherian $k$-基底への強 Noether 性や、任意 AS/QP 表示との
同一視は、この系から追加で主張しない。
\end{proof}
